\section{Introduction}

Large language models have achieved remarkable success under the autoregressive paradigm \citep{brown2020language,guo2025deepseek,liu2024deepseekv2,hurst2024gpt,team2026kimi}. By factorizing the discrete text distribution through the chain rule \citep{malach2023auto,barrault2024large,katharopoulos2020transformers,hutchins2022block,du2022glm,yang2019xlnet}, autoregressive language models have driven major advances in large-scale pretraining, open-ended generation, and downstream transfer, and have become the dominant approach to modern language modeling \citep{pan2025survey,minaee2024large,chang2024survey,wan2023efficient,zhao2023survey}. However, this paradigm tightly couples generation to a fixed left-to-right order, making inference inherently sequential and restricting the model’s inductive bias to a single token ordering \citep{bachmann2024pitfalls,wu2025constrained,li2025beyond,monea2023pass,aishwarya2024tandem,fu2024break,you2024linear}. Recent progress in both discrete and continuous diffusion-based text modeling suggests that high-quality language generation need not rely on such a fixed order; instead, language models can also be defined through more general state evolution and denoising paths \citep{ye2024beyond,yu2025dimple,ou2024your,strudel2022self,chen2025dlm,mahabadi2024tess}.

Despite extensive exploration along autoregressive, discrete diffusion, and continuous diffusion directions \citep{song2020score,von2025generalized,li2025unifying,xu2024energy,goel2026skip,jarrett2025language,jeoung2025examining,lahoti2023improving}, existing methods still struggle to simultaneously achieve generation efficiency, scalable representation, and global semantic modeling. Autoregressive models directly parameterize token-level conditional probabilities, yielding a clear training objective, but their fixed generation order incurs inherent sequential inference cost and introduces a strong hand-crafted inductive bias, which limits performance on more general generation tasks \citep{lin2021limitations,berglund2023reversal,deschenaux2024promises,dalal2019autoregressive,zhu2024towards}. Discrete diffusion language models remove explicit left-to-right factorization \citep{hoogeboom2021autoregressive,gat2024discrete,hoogeboom2021argmax,zhao2025d1}, yet they still typically perform observation recovery in discrete token space, leading to costly multi-step sampling, while intermediate discrete states are not well suited to stably represent global semantic structure \citep{zheng2023reparameterized,zhou2024diffusion,venkatraman2024amortizing,srikanth2025discrete,takida2022sq,meng2022concrete,jang2016categorical}. Continuous diffusion methods further introduce continuous representation spaces \citep{sahoo2025diffusion,gong2024scaling,tae2025tess}, but most existing approaches still use the diffusion path to recover token-aligned representations rather than to explicitly model a latent prior \citep{gulrajani2023likelihood,dieleman2022continuous}. As a result, current methods have not yet provided a unified framework that systematically combines non-autoregressive generation, continuous representation, and probabilistic text modeling.

To address this gap, we propose \method, a hierarchical latent-space diffusion language model. \method first learns a stable mapping between text and continuous latent variables through a Text VAE \citep{zheng2025diffusion,xu2020variational,semeniuta2017hybrid,li2020optimus,bowman2016generating,kingma2013auto}, then models the latent prior in continuous latent space with a block-causal DiT \citep{peebles2023scalable,chandrasegaran2025exploring,park2024switch,mo2023dit,liu2025bcat,cannizzaro2023towards,bandyopadhyay2025block,zhang2020invariant}, and finally generates text through a conditional decoder. The key idea of \method is to use diffusion not for token-level observation recovery, but for latent prior transport. From a unified Markov-path perspective, this design explicitly decomposes text generation into two levels: global semantic organization in continuous latent space and local textual realization through conditional decoding. This decomposition weakens the inductive bias imposed by fixed token order, allows the geometry of continuous space to directly support semantic compression and prior fitting, and enables a more flexible non-autoregressive generation process. Moreover, block-causal prior modeling preserves cross-block causal structure while allowing more efficient parallel computation within each block. Grounded in hierarchical latent-space modeling, \method is also highly modular and naturally extensible to alternative latent modeling components and other continuous modalities \citep{zheng2025diffusion,deng2026generative}.

Motivated by these observations, we systematically study diffusion language modeling in continuous latent space from both theoretical and empirical perspectives. Our contributions are as follows.

\begin{itemize}
     \item We propose \method, a hierarchical latent-space language model that explicitly decomposes text generation into global semantic modeling and local textual realization within a unified probabilistic framework, while using diffusion-based prior modeling in continuous latent space to connect the two, thereby establishing a new paradigm for language generation from the perspective of hierarchical information decomposition.
    \item We analyze the differences between \method and existing language modeling paradigms from a unified Markov-path perspective, clarifying its advantages in global semantic modeling, non-autoregressive inductive bias, and theoretical interpretability, which are further validated in the subsequent experiments.
    \item Through extensive experiments spanning 4 research questions, 8 benchmarks, strictly matched $\sim$2B-parameter autoregressive and LLaDA baselines, and scaling curves up to about 2000 EFLOPs, we systematically validate the central claims of \method, identify an effective overall configuration, and verify its strong potential and favorable scaling behavior for text generation.
    \item We further analyze several issues beyond the core framework, including the mismatch between likelihood estimation and generation quality, first-block conditioning, and latent compression. We also provide preliminary evidence that \method offers a natural bridge from discrete text to continuous modalities such as vision, pointing to a broader unified generative paradigm.
\end{itemize}